% ─── Capítulo: Manual de Usuario ──────────────────────────────────────────────
\chapter{MANUAL DE USUARIO}

\section{Sintaxis del Lenguaje TPP}

El lenguaje TPP utiliza palabras reservadas en español y sigue una estructura similar a C o Java. Todo archivo fuente debe tener la extensión \texttt{.to}. A continuación se detalla cómo utilizar cada una de las estructuras del lenguaje.

\subsection{Tipos de Datos y Variables}

TPP soporta variables \texttt{entero} y \texttt{decimal}, así como arreglos estáticos. Todas las sentencias deben terminar en punto y coma (\texttt{;}).

\begin{lstlisting}[style=tppstyle, caption={Declaración de variables y arreglos}]
// Declaracion e inicializacion
entero edad = 20;
decimal pi = 3.14;

// Declaracion multiple
entero x, y, z;

// Declaracion de arreglos (el tamano debe ser un numero literal)
entero notas[5];
notas[0] = 15;
notas[1] = 18;
\end{lstlisting}

\subsection{Entrada y Salida}

\begin{lstlisting}[style=tppstyle, caption={Operaciones de Entrada y Salida}]
// Imprimir cadenas de texto o variables
imprimir("Bienvenido al sistema");
imprimir(edad);

// Leer un valor desde consola hacia una variable
leer(edad);
\end{lstlisting}

\subsection{Expresiones y Operadores}

El lenguaje TPP permite construir expresiones aritméticas, lógicas y relacionales complejas respetando la precedencia matemática estándar (de mayor a menor prioridad):

\begin{itemize}
    \item \textbf{Agrupación:} Paréntesis \texttt{( )} para forzar la prioridad de evaluación.
    \item \textbf{Aritméticos:} Multiplicación \texttt{*}, División \texttt{/} (prioridad alta). Suma \texttt{+}, Resta \texttt{-} (prioridad baja).
    \item \textbf{Relacionales:} Mayor \texttt{>}, Menor \texttt{<}, Mayor o Igual \texttt{>=}, Menor o Igual \texttt{<=}.
    \item \textbf{Igualdad:} Igual que \texttt{==}, Diferente de \texttt{!=}.
\end{itemize}

\begin{infobox}[Tipos en Expresiones]
TPP realiza \textbf{conversiones implícitas (coerción)} de tipos cuando se mezclan enteros y decimales en una operación aritmética. El resultado automáticamente se promoverá a decimal para no perder precisión. Sin embargo, los arreglos no pueden ser operados directamente (solo sus elementos individuales).
\end{infobox}

\subsection{Estructuras Condicionales: \texttt{si} / \texttt{sino}}

La estructura condicional evalúa una expresión. Si es diferente de cero (verdadero), ejecuta el primer bloque; de lo contrario, el bloque \texttt{sino}.

\begin{lstlisting}[style=tppstyle, caption={Uso de si y sino}]
si (edad >= 18) {
    imprimir("Es mayor de edad");
} sino {
    imprimir("Es menor de edad");
}
\end{lstlisting}

\subsection{Selección Múltiple: \texttt{depende} / \texttt{caso}}

Equivalente a la construcción \texttt{switch/case} de lenguajes de la familia C. Permite evaluar el valor de una variable y ejecutar exclusivamente el bloque que coincida con el caso correspondiente.

\textbf{Detalles importantes:}
\begin{itemize}
    \item A diferencia de C, \textbf{no es necesario utilizar un equivalente a \texttt{break}}. La semántica de TPP asume que, una vez que un caso se ejecuta, el flujo del programa sale automáticamente del bloque \texttt{depende}.
    \item Los casos solo pueden comparar contra literales constantes, no contra otras variables.
    \item El bloque \texttt{otros} actúa de forma idéntica a \texttt{default}, ejecutándose si ninguna de las condiciones anteriores se cumplió. Su presencia es opcional.
\end{itemize}

\begin{lstlisting}[style=tppstyle, caption={Uso de depende y caso}]
entero opcion = 2;
depende (opcion) {
    caso 1: {
        imprimir("Opcion 1 seleccionada");
    }
    caso 2: {
        imprimir("Opcion 2 seleccionada");
    }
    otros: {
        imprimir("Opcion no valida");
    }
}
\end{lstlisting}

\subsection{Bucles: \texttt{mientras} y \texttt{para}}

Los bucles permiten repetir un bloque de código. TPP soporta dos construcciones iterativas principales.

\textbf{1. Bucle \texttt{mientras} (while):} 
Se evalúa la condición de continuación antes de cada iteración. Si es verdadera, se ejecuta el cuerpo del bucle. Útil cuando no se conoce de antemano el número de iteraciones.

\textbf{2. Bucle \texttt{para} (for):} 
Es un bloque compacto que concentra la lógica de iteración. Requiere exactamente tres partes separadas por punto y coma:
\begin{enumerate}
    \item \textbf{Inicialización:} Se ejecuta una sola vez al entrar al bucle. Puede declarar una nueva variable iteradora (ej. \texttt{entero i = 0}) o inicializar una existente.
    \item \textbf{Condición:} Expresión booleana evaluada antes de cada iteración.
    \item \textbf{Actualización:} Instrucción que se ejecuta al final de cada iteración, típicamente para incrementar el contador.
\end{enumerate}

\begin{lstlisting}[style=tppstyle, caption={Uso de bucles}]
// Bucle MIENTRAS
entero contador = 0;
mientras (contador < 5) {
    imprimir(contador);
    contador = contador + 1; // No existe ++ en TPP
}

// Bucle PARA (inicializacion; condicion; incremento)
para (entero i = 0; i < 10; i = i + 1) {
    imprimir(i);
}
\end{lstlisting}

\subsection{Funciones: \texttt{fun}}

Las funciones se declaran con la palabra \texttt{fun}. Pueden tener parámetros y opcionalmente devolver un valor con la flecha \texttt{->}. La palabra \texttt{retornar} devuelve el resultado.

\begin{lstlisting}[style=tppstyle, caption={Definicion de funciones}]
// Funcion que retorna un valor
fun sumar(entero a, entero b) -> entero {
    retornar a + b;
}

// Funcion sin valor de retorno (procedimiento)
fun saludar(entero veces) {
    para (entero i = 0; i < veces; i = i + 1) {
        imprimir("Hola");
    }
}

// Punto de entrada principal (opcional, por convención)
fun main() {
    entero res = sumar(5, 3);
    imprimir(res);
}
\end{lstlisting}

\section{Guía de Compilación y Ejecución en Logisim}

Para llevar un programa fuente \texttt{.to} hasta el procesador simulado en Logisim, siga estos pasos:

\subsection{1. Compilación del código fuente}

Ejecute el compilador pasándole el archivo \texttt{.to}. Si es exitoso, generará el archivo ensamblador \texttt{output.s}.

\begin{lstlisting}[language=bash, style=cstyle]
$ ./tpp programa.to
\end{lstlisting}
Si hay errores léxicos, sintácticos o semánticos, el compilador los reportará en la consola y abortará el proceso.

\subsection{2. Generación del Código Hexadecimal}

El ensamblador de Logisim requiere que las instrucciones estén en formato hexadecimal (\texttt{v2.0 raw}). Se utiliza el script de Python incluido en el repositorio:

\begin{lstlisting}[language=bash, style=cstyle]
$ python3 compile_to_logisim.py
\end{lstlisting}

El script leerá \texttt{output.s} y generará archivos hexadecimales (por defecto produce \texttt{output.hex}, o \texttt{rom.hex} y \texttt{ram.hex} si la arquitectura del circuito separa las memorias de instrucciones y datos).

\subsection{3. Carga en Logisim Evolution}

\begin{enumerate}
    \item Abra su proyecto de procesador RV32I en \textbf{Logisim Evolution}.
    \item Localice el componente de memoria \textbf{ROM} (memoria de instrucciones) y/o \textbf{RAM} (memoria de datos).
    \item Haga \textbf{clic derecho} sobre el componente de memoria.
    \item Seleccione \textbf{Load Image...} (Cargar Imagen).
    \item Seleccione el archivo generado (\texttt{rom.hex}, \texttt{ram.hex} u \texttt{output.hex} según corresponda).
    \item Inicie el reloj (Ctrl+K o menú Simulate $\rightarrow$ Auto-Tick) para ver el programa ejecutarse.
    \item Si el programa utiliza la instrucción \texttt{imprimir}, abra el componente TTY para ver la salida del programa.
\end{enumerate}
